\documentclass[12pt,a4paper]{article}

% ─── Packages ───────────────────────────────────────────────────────────────
\usepackage[a4paper, margin=1in]{geometry}
\usepackage{amsmath, amssymb}
\usepackage{graphicx}
\usepackage{booktabs}
\usepackage{array}
\usepackage{tabularx}
\usepackage{multirow}
\usepackage{xcolor}
\usepackage{hyperref}
\usepackage{caption}
\usepackage{subcaption}
\usepackage{fancyhdr}
\usepackage{titlesec}
\usepackage{setspace}
\usepackage{enumitem}
\usepackage{float}
\usepackage{siunitx}

% ─── Colours ────────────────────────────────────────────────────────────────
\definecolor{headerblue}{RGB}{31,73,125}
\definecolor{rowgray}{RGB}{235,235,235}

% ─── Header / Footer ────────────────────────────────────────────────────────
\pagestyle{fancy}
\fancyhf{}
\rhead{\textcolor{headerblue}{\small Project 3 -- Constant $g_m$ Bias}}
\lhead{\textcolor{headerblue}{\small EE Department, IIT Kanpur}}
\cfoot{\thepage}
\renewcommand{\headrulewidth}{0.4pt}

% ─── Section Formatting ──────────────────────────────────────────────────────
\titleformat{\section}{\large\bfseries\color{headerblue}}{}{0em}{}[\titlerule]
\titleformat{\subsection}{\normalsize\bfseries}{}{0em}{}

% ─── Hyperlinks ──────────────────────────────────────────────────────────────
\hypersetup{colorlinks=true, linkcolor=headerblue, urlcolor=headerblue}

% ─── Spacing ─────────────────────────────────────────────────────────────────
\setlength{\parskip}{6pt}
\setlength{\parindent}{0pt}
\onehalfspacing

% ════════════════════════════════════════════════════════════════════════════
\begin{document}

% ─── Title Page ──────────────────────────────────────────────────────────────
\begin{titlepage}
  \centering
  \vspace*{1.5cm}

  {\Huge\bfseries\color{headerblue} Project Report}\\[0.5cm]
  {\LARGE\bfseries Project 3: Constant $g_m$ Bias}\\[0.3cm]
  {\large 4-Transistor OTA via Negative Feedback}\\[1.2cm]

  \rule{\linewidth}{0.5pt}\\[0.5cm]

  \begin{tabular}{ll}
    \textbf{Course}      & Analog Integrated Circuit Design \\[4pt]
    \textbf{Department}  & Electrical Engineering \\[4pt]
    \textbf{Institute}   & Indian Institute of Technology Kanpur \\[4pt]
    \textbf{Tool}        & Cadence Virtuoso (Spectre Simulator) \\[4pt]
    \textbf{Technology}  & CMOS (1.8\,V supply) \\
  \end{tabular}

  \vspace{1cm}
  \rule{\linewidth}{0.5pt}

  \vfill
  {\large \today}
\end{titlepage}

% ─── Table of Contents ───────────────────────────────────────────────────────
\tableofcontents
\newpage

% ════════════════════════════════════════════════════════════════════════════
\section{Introduction}

Transconductance ($g_m$) is one of the most critical parameters in analog circuit design. It
directly governs the gain, bandwidth, and noise performance of amplifiers, filters, and
data-converters. However, $g_m$ is inherently sensitive to process, voltage, and temperature
(PVT) variations. A transistor biased with a fixed current will show significant $g_m$ drift
because $g_m = \sqrt{2\mu_n C_{ox}(W/L)I_D}$ for an NMOS in saturation, where mobility
$\mu_n$ decreases with temperature.

This project addresses the problem by designing a \textbf{constant-$g_m$ bias circuit} using a
compact \textbf{4-transistor OTA (VCCS)} architecture governed by a negative feedback servo loop.
The feedback loop continuously senses the $g_m$ of a replica OTA and adjusts the bias such that

\begin{equation}
  G_m = \frac{I_{\text{ref}}}{2\,\Delta V}
\end{equation}

where $I_{\text{ref}}$ is a stable current reference and $2\Delta V$ is a fixed incremental
voltage generated from a resistive voltage divider referenced to $V_{\text{DD}}$.

\subsection*{Design Specifications}

\begin{table}[H]
  \centering
  \caption{Project Performance Targets}
  \renewcommand{\arraystretch}{1.3}
  \begin{tabular}{lll}
    \toprule
    \textbf{Parameter}        & \textbf{Target}              & \textbf{Constraint} \\
    \midrule
    Architecture              & 4-transistor OTA (VCCS)      & --                  \\
    Target $g_m$              & $10\,\text{mS}$              & At nominal PVT      \\
    Control Mechanism         & Circuit-based negative FB    & Self-regulating     \\
    $g_m$ Variation           & $< 1\%$                      & $-20\,^\circ$C to $+100\,^\circ$C \\
    Supply Voltage ($V_{DD}$) & $1.8\,\text{V}$              & --                  \\
    \bottomrule
  \end{tabular}
\end{table}

% ════════════════════════════════════════════════════════════════════════════
\section{Working Principle}

\subsection{Negative Feedback Servo Loop}

The core idea is borrowed from a well-known servo-loop technique for $g_m$ setting. A small
incremental DC voltage $2\Delta V$ is applied differentially to the input of a replica 4T OTA.
The resulting output current $2G_m\Delta V$ is compared against a stable reference current
$I_{\text{ref}}$ using an error amplifier $G_{\text{err}}$.

The error amplifier drives the common-mode bias voltage $V_{cm}$ of the OTA. If $G_m$
increases (due to temperature drop or process variation), the output current exceeds
$I_{\text{ref}}$, causing $V_{cm}$ to decrease, which in turn reduces $G_m$ — a classic
negative feedback action. At steady state:

\begin{equation}
  2\,G_m\,\Delta V = I_{\text{ref}}
  \quad\Longrightarrow\quad
  G_m = \frac{I_{\text{ref}}}{2\,\Delta V}
\end{equation}

Since $I_{\text{ref}}$ and $\Delta V$ are both derived from stable references, $G_m$ is
locked to a PVT-insensitive value.

\subsection{Voltage Divider for $2\Delta V$}

The incremental voltage is generated from a symmetric resistive divider:

\begin{equation}
  2\Delta V = V_{\text{ref}}\,\frac{R_1}{R + R_1}
\end{equation}

In the implementation, $R_0 = R_3 = 8.989\,\text{k}\Omega$ and $R_1 = R_2 = 100\,\Omega$,
giving a small $\Delta V$ that keeps the OTA well within its linear region.

\subsection{Common-Mode Feedback}

A common-mode feedback (CMFB) path sets the output common-mode voltage $V_{op,om}$ to $V_{cm}$.
The two op-amp blocks (I4 and I5) serve as error amplifiers: one ($G_{err}$) drives the $g_m$
servo loop, while the other enforces the output common-mode level via a feedback path.

% ════════════════════════════════════════════════════════════════════════════
\section{Circuit Architecture}

\subsection{Top-Level Negative Feedback Loop (negative\_fb\_loop)}

The top-level schematic implements the complete servo system:

\begin{itemize}[leftmargin=1.5em]
  \item \textbf{BLOCK} -- The 4-transistor OTA replica whose $g_m$ is being regulated.
  \item \textbf{R0, R3} ($8.989\,\text{k}\Omega$) -- Upper/lower arms of the voltage divider.
  \item \textbf{R1, R2} ($100\,\Omega$) -- Middle taps that define $\pm\Delta V$.
  \item \textbf{R4, R5} ($1\,\text{M}\Omega$) -- Averaging resistors for $V_{op}$ and $V_{om}$.
  \item \textbf{I1, I2} ($94\,\mu\text{A}$) -- Reference current sources ($I_{\text{ref}}$).
  \item \textbf{I4, I5} -- Error amplifiers (Opamp instances) closing the servo and CMFB loops.
  \item \textbf{C1} ($2\,\text{pF}$) -- Compensation capacitor for loop stability.
\end{itemize}

\subsection{4-Transistor OTA (BLOCK)}

The BLOCK is a differential-pair OTA with a current-mirror load:

\begin{itemize}[leftmargin=1.5em]
  \item \textbf{PM0, PM1} (PMOS, $W=6\,\mu$m, $L=540\,\text{nm}$, $m=24$) -- Current mirror load.
  \item \textbf{NM0, NM1} (NMOS, $W=9\,\mu$m, $L=540\,\text{nm}$, $m=14$) -- Differential pair.
  \item \textbf{R0} ($100\,\Omega$) -- Source degeneration resistor to linearise $G_m$.
  \item Ports: \texttt{Vin\_Plus}, \texttt{Vin\_minus}, \texttt{Vcm\_fb}, \texttt{Vom}, \texttt{Vop}.
\end{itemize}

\subsection{Error Amplifier (Opamp -- I4 / I5)}

Each error amplifier is a standard 5-transistor OTA:

\begin{itemize}[leftmargin=1.5em]
  \item \textbf{PM0, PM1} (PMOS, $W=2\,\mu$m, $L=2\,\mu$m, $m=11$) -- Current mirror load.
  \item \textbf{NM0, NM1} (NMOS, $W=2\,\mu$m, $L=2\,\mu$m, $m=11$) -- Differential pair.
  \item \textbf{I0} ($20\,\mu\text{A}$) -- Tail current source.
  \item Ports: \texttt{Vin\_plus}, \texttt{Vin\_minus}, \texttt{Vout}.
\end{itemize}

% ════════════════════════════════════════════════════════════════════════════
\section{Operating Point Analysis}

All operating points were extracted using Cadence Spectre DC analysis at nominal temperature
($27\,^\circ$C) and $V_{DD} = 1.8\,\text{V}$.

\subsection{4T OTA (BLOCK) -- Transistor Operating Points}

\begin{table}[H]
  \centering
  \caption{DC Operating Points -- 4T OTA (BLOCK)}
  \renewcommand{\arraystretch}{1.4}
  \begin{tabularx}{\textwidth}{lXXXX}
    \toprule
    \rowcolor{headerblue}
    \textcolor{white}{\textbf{Parameter}}
      & \textcolor{white}{\textbf{PM0 (PMOS)}}
      & \textcolor{white}{\textbf{PM1 (PMOS)}}
      & \textcolor{white}{\textbf{NM0 (NMOS)}}
      & \textcolor{white}{\textbf{NM1 (NMOS)}} \\
    \midrule
    \rowcolor{rowgray}
    Role              & Mirror Load & Mirror Load & Diff. Pair & Diff. Pair \\
    Region            & Saturation  & Saturation  & Saturation & Saturation \\
    \rowcolor{rowgray}
    $V_{GS}$ (mV)    & $-831.1$    & $-831.1$    & $+758.9$   & $+739.1$   \\
    $V_{DS}$ (mV)    & $-796.9$    & $-884.1$    & $+792.5$   & $+705.3$   \\
    \rowcolor{rowgray}
    $V_{TH}$ (mV)    & $-464.2$    & $-464.2$    & $+560.0$   & $+560.0$   \\
    $V_{DSAT}$ (mV)  & $-324.1$    & $-324.1$    & $+204.8$   & $+190.4$   \\
    \rowcolor{rowgray}
    $I_{DS}$ ($\mu$A) & $-1050.1$ & $-1056.2$   & $+1144.0$  & $+962.2$   \\
    $g_m$ (mS)       & $5.028$     & $5.058$     & $9.249$    & $8.480$    \\
    \rowcolor{rowgray}
    $g_{ds}$ ($\mu$S) & $73.38$   & $66.43$     & $73.30$    & $69.94$    \\
    $g_m/I_D$        & $4.789$     & $4.789$     & $8.084$    & $8.813$    \\
    \rowcolor{rowgray}
    $r_{out}$ (k$\Omega$) & $13.63$ & $15.05$  & $13.64$    & $14.30$    \\
    Self Gain ($g_m r_{out}$) & $68.5$ & $76.1$ & $126.2$   & $121.2$    \\
    \bottomrule
  \end{tabularx}
\end{table}

\textbf{Observations:}
\begin{itemize}[leftmargin=1.5em]
  \item All four transistors are firmly in \textbf{saturation}, as required for proper OTA operation.
  \item The NMOS differential pair (NM0/NM1) exhibits higher $g_m$ ($\approx 8.5$--$9.2\,\text{mS}$)
        than the PMOS mirror load ($\approx 5.0\,\text{mS}$), consistent with the
        higher electron mobility.
  \item The $g_m/I_D$ ratio for NMOS ($\approx 8$--$8.8$) indicates moderate inversion, while
        PMOS ($\approx 4.8$) operates at a higher overdrive, which is expected for a current
        mirror with wide transistors.
  \item Source degeneration ($R_0 = 100\,\Omega$) linearises the overall $G_m$ at the cost of
        a slight reduction in effective transconductance.
\end{itemize}

% ──────────────────────────────────────────────────────────────────────────
\subsection{Error Amplifier I5 -- Transistor Operating Points}

\begin{table}[H]
  \centering
  \caption{DC Operating Points -- Error Amplifier I5 (Opamp)}
  \renewcommand{\arraystretch}{1.4}
  \begin{tabularx}{\textwidth}{lXXXX}
    \toprule
    \rowcolor{headerblue}
    \textcolor{white}{\textbf{Parameter}}
      & \textcolor{white}{\textbf{PM0 (PMOS)}}
      & \textcolor{white}{\textbf{PM1 (PMOS)}}
      & \textcolor{white}{\textbf{NM0 (NMOS)}}
      & \textcolor{white}{\textbf{NM1 (NMOS)}} \\
    \midrule
    \rowcolor{rowgray}
    Role              & Mirror Load & Mirror Load & Diff. Pair & Diff. Pair \\
    Region            & Saturation  & Saturation  & Saturation & Saturation \\
    \rowcolor{rowgray}
    $V_{GS}$ (mV)    & $-598.6$    & $-598.6$    & $+582.4$   & $+583.3$   \\
    $V_{DS}$ (mV)    & $-598.6$    & $-831.1$    & $+824.3$   & $+591.8$   \\
    \rowcolor{rowgray}
    $V_{TH}$ (mV)    & $-437.3$    & $-437.2$    & $+559.4$   & $+559.5$   \\
    $V_{DSAT}$ (mV)  & $-167.8$    & $-167.8$    & $+97.4$    & $+97.9$    \\
    \rowcolor{rowgray}
    $I_{DS}$ ($\mu$A) & $-9.97$   & $-10.03$    & $+9.97$    & $+10.03$   \\
    $g_m$ ($\mu$S)   & $98.81$     & $99.30$     & $174.3$    & $175.1$    \\
    \rowcolor{rowgray}
    $g_{ds}$ (nS)    & $281.2$     & $247.2$     & $351.9$    & $386.3$    \\
    $g_m/I_D$        & $9.91$      & $9.90$      & $17.48$    & $17.46$    \\
    \rowcolor{rowgray}
    $r_{out}$ (M$\Omega$) & $3.56$ & $4.05$   & $2.84$     & $2.59$     \\
    Self Gain         & $351.4$     & $401.7$     & $495.3$    & $453.3$    \\
    \bottomrule
  \end{tabularx}
\end{table}

% ──────────────────────────────────────────────────────────────────────────
\subsection{Error Amplifier I4 -- Transistor Operating Points}

\begin{table}[H]
  \centering
  \caption{DC Operating Points -- Error Amplifier I4 (Opamp)}
  \renewcommand{\arraystretch}{1.4}
  \begin{tabularx}{\textwidth}{lXXXX}
    \toprule
    \rowcolor{headerblue}
    \textcolor{white}{\textbf{Parameter}}
      & \textcolor{white}{\textbf{PM0 (PMOS)}}
      & \textcolor{white}{\textbf{PM1 (PMOS)}}
      & \textcolor{white}{\textbf{NM0 (NMOS)}}
      & \textcolor{white}{\textbf{NM1 (NMOS)}} \\
    \midrule
    \rowcolor{rowgray}
    Role              & Mirror Load & Mirror Load & Diff. Pair  & Output/Tail \\
    Region            & Saturation  & Saturation  & Saturation  & Saturation* \\
    \rowcolor{rowgray}
    $V_{GS}$ (mV)    & $-656.4$    & $-656.4$    & $+615.8$    & $+528.7$   \\
    $V_{DS}$ (mV)    & $-656.4$    & $-839.7$    & $+756.3$    & $+573.1$   \\
    \rowcolor{rowgray}
    $V_{TH}$ (mV)    & $-437.2$    & $-437.2$    & $+561.9$    & $+562.0$   \\
    $V_{DSAT}$ (mV)  & $-211.0$    & $-211.1$    & $+116.7$    & $+71.6$    \\
    \rowcolor{rowgray}
    $I_{DS}$ ($\mu$A) & $-16.60$  & $-16.67$    & $+16.60$    & $+3.40$    \\
    $g_m$ ($\mu$S)   & $129.7$     & $130.2$     & $259.9$     & $68.9$     \\
    \rowcolor{rowgray}
    $g_{ds}$ (nS)    & $414.3$     & $364.7$     & $534.8$     & $155.7$    \\
    $g_m/I_D$        & $7.81$      & $7.81$      & $15.66$     & $20.27$    \\
    \rowcolor{rowgray}
    $r_{out}$ (M$\Omega$) & $2.41$ & $2.74$   & $1.87$      & $6.42$     \\
    Self Gain         & $313.1$     & $356.9$     & $486.0$     & $442.4$    \\
    \bottomrule
  \end{tabularx}
  \begin{flushleft}
    \small *NM1 in I4 operates at region code 3 (near saturation/subthreshold boundary)
    due to its low $I_{DS} = 3.4\,\mu$A compared to the main pair.
  \end{flushleft}
\end{table}

\textbf{Observations on Amplifiers I4 and I5:}
\begin{itemize}[leftmargin=1.5em]
  \item Both amplifiers are biased at low currents ($\approx 10$--$17\,\mu$A per branch),
        resulting in high output resistances in the M$\Omega$ range, which translates to
        large DC gain well suited for the feedback loop.
  \item The very high self-gain values (300--495) indicate that both op-amps provide
        sufficient loop gain to enforce $G_m = I_{\text{ref}}/2\Delta V$ accurately.
  \item NMOS transistors in both amplifiers operate in moderate inversion
        ($g_m/I_D \approx 15$--$20$), prioritising gain efficiency over speed.
  \item $V_{DSAT}$ values are small ($< 220\,\text{mV}$), ensuring adequate output
        swing headroom within the $1.8\,\text{V}$ supply.
\end{itemize}

% ════════════════════════════════════════════════════════════════════════════
\section{Simulation Results and Specification Status}

\subsection{$g_m$ vs. Temperature — Observed Result}

A DC sweep simulation was performed over temperature from $-20\,^\circ$C to $+100\,^\circ$C
using Cadence Spectre. The transconductance was probed directly from the device operating
point of NM0 (\texttt{OS("/I0/NM0" "gm")}) in the closed-loop \texttt{negative\_fb\_loop}
schematic.

\begin{table}[H]
  \centering
  \caption{Simulated $g_m$ vs.\ temperature — actual results}
  \renewcommand{\arraystretch}{1.3}
  \begin{tabular}{cccc}
    \toprule
    \textbf{Temp ($^\circ$C)} & \textbf{$g_m$ (mS)} & \textbf{Deviation from 27$^\circ$C} & \textbf{Spec ($<1\%$)} \\
    \midrule
    $-20$  & $9.80$ & $+5.95\%$  & \textcolor{red}{\textbf{FAIL}} \\
    $0$    & $9.55$ & $+3.24\%$  & \textcolor{red}{\textbf{FAIL}} \\
    $27$   & $9.25$ & Reference  & --                              \\
    $50$   & $8.90$ & $-3.78\%$  & \textcolor{red}{\textbf{FAIL}} \\
    $75$   & $8.55$ & $-7.57\%$  & \textcolor{red}{\textbf{FAIL}} \\
    $100$  & $8.20$ & $-11.35\%$ & \textcolor{red}{\textbf{FAIL}} \\
    \bottomrule
  \end{tabular}
\end{table}

\textbf{Result:} The $g_m$ varies from $\approx 9.80\,\text{mS}$ at $-20\,^\circ$C to
$\approx 8.20\,\text{mS}$ at $+100\,^\circ$C — a total variation of $\mathbf{\sim 19\%}$,
far exceeding the target of $< 1\%$. The design \textbf{does not meet} the specification.

\subsection{Discussion — Why the $< 1\%$ Specification is Not Met}

The dominant reason the specification is not achieved is the \textbf{temperature dependence
of carrier mobility}. In a MOSFET biased at a fixed drain current, the transconductance is:
\[
  g_m = \sqrt{2\,\mu_n(T)\,C_{ox}\,(W/L)\,I_{DS}}
\]
where $\mu_n(T) \propto T^{-3/2}$ decreases with rising temperature. Since the bias current
$I_{DS}$ in the 4T OTA is set by the tail source — which itself has a finite temperature
coefficient in a real implementation — $g_m$ falls nearly linearly with temperature, producing
the observed $\sim 19\%$ drift.

The negative feedback servo loop is intended to counteract this by increasing $I_{DS}$
(via raising $V_{cm\_fb}$) as temperature rises, so that the product
$\mu_n(T) \cdot I_{DS}(T)$ remains constant. However, achieving this requires the loop to
track the mobility variation \emph{precisely} and in \emph{real time}, which places a
stringent demand on the loop gain and the accuracy of $I_{\text{ref}}$ and $\Delta V$.

In this implementation, the available loop correction is insufficient to fully compensate
the steep mobility roll-off across the $-20\,^\circ$C to $+100\,^\circ$C range. Additionally,
polysilicon resistors used to generate $\Delta V$ carry a non-negligible temperature
coefficient ($\approx +1500\,\text{ppm/}^\circ$C), which causes $\Delta V$ — and hence
the target locking point $G_m = I_{\text{ref}}/2\Delta V$ — to shift slightly with
temperature, further degrading the regulation accuracy. These two effects together account
for the observed gap between the simulated performance and the $< 1\%$ target.

\subsection{Common-Mode Voltages (net6/net8) vs. Temperature}

The feedback node voltages (net6, net8) vary in the range $965$--$980\,\text{mV}$ across
temperature, showing a non-monotonic U-shaped profile. Ideally, if the servo loop were
functioning correctly, $V_{cm\_fb}$ should sweep significantly (by several hundred mV) to
actively compensate the falling $g_m$. The narrow variation observed ($\sim 15\,\text{mV}$)
strongly suggests the loop is not exerting meaningful control over the bias point.



% ════════════════════════════════════════════════════════════════════════════
\section{Design Calculations}

\subsection{Correct $g_m$ Locking Point}

For the servo to lock $G_m = 10\,\text{mS}$, the required $2\Delta V$ with
$I_{\text{ref}} = 94\,\mu\text{A}$ is:
\begin{equation}
  2\Delta V_{\text{required}} = \frac{I_{\text{ref}}}{G_m}
  = \frac{94\,\mu\text{A}}{10\,\text{mS}} = 9.4\,\text{mV}
\end{equation}

The current resistor values produce:
\begin{equation}
  2\Delta V_{\text{actual}} = 1.8\,\text{V} \times \frac{100}{9089} \approx 19.8\,\text{mV}
  \quad\Rightarrow\quad G_{m,\text{locked}} \approx 4.75\,\text{mS}
\end{equation}

To correct this without changing $I_{\text{ref}}$, $R_1$ should be reduced:
\begin{equation}
  R_1 = \frac{2\Delta V_{\text{required}}}{V_{DD} - 2\Delta V_{\text{required}}}
        \times R
  = \frac{9.4\,\text{mV}}{1.8\,\text{V} - 9.4\,\text{mV}} \times 8989\,\Omega
  \approx 47\,\Omega
\end{equation}

Alternatively, $I_{\text{ref}}$ can be increased to $198\,\mu\text{A}$ keeping the
existing resistor values.

\subsection{Effective $G_m$ with Source Degeneration}

With source degeneration resistor $R_S = 100\,\Omega$ (split: $50\,\Omega$ per side),
the effective OTA transconductance is:
\begin{equation}
  G_{m,\text{eff}} = \frac{g_m}{1 + g_m\,(R_S/2)}
  = \frac{9.25\,\text{mS}}{1 + 9.25\,\text{mS} \times 50\,\Omega}
  \approx 8.13\,\text{mS}
\end{equation}

The servo loop must account for this degeneration when locking $G_m$ to the target.

\subsection{Output Resistance of 4T OTA}

\begin{equation}
  R_{\text{out}} = r_{o,\text{NM1}} \,\|\, r_{o,\text{PM1}}
  = 14.30\,\text{k}\Omega \,\|\, 15.05\,\text{k}\Omega
  \approx 7.33\,\text{k}\Omega
\end{equation}

\subsection{Required Loop Gain for $< 1\%$ Regulation}

To suppress the inherent $19\%$ open-loop variation down to $< 1\%$:
\begin{equation}
  \frac{\Delta G_m}{G_m}\bigg|_{\text{CL}} =
  \frac{\Delta G_m / G_m|_{\text{OL}}}{1 + T}
  < 1\%
  \quad\Rightarrow\quad
  1 + T > \frac{19\%}{1\%} = 19
  \quad\Rightarrow\quad T > 18
\end{equation}

With self-gain of NM0 $\approx 126$ and the error amplifier gain $> 300$, the available
loop gain is theoretically sufficient. The issue lies in how effectively this gain is
utilised in the actual feedback path.



% ════════════════════════════════════════════════════════════════════════════
\section{Summary and Conclusion}

This project involved the design and simulation of a constant-$g_m$ bias circuit using a
compact 4-transistor OTA with a negative feedback servo loop in Cadence Virtuoso Spectre.

\subsection*{What Was Achieved}
\begin{itemize}[leftmargin=1.5em]
  \item A 4T OTA with NMOS differential pair and PMOS mirror load was successfully designed
        with all transistors biased in \textbf{saturation} at nominal conditions
        ($V_{DD} = 1.8\,\text{V}$, $27\,^\circ$C).
  \item Two high-gain error amplifiers (I4, I5) with individual transistor self-gain
        $> 300$ were designed and verified through operating point extraction.
  \item The complete servo loop schematic — including voltage divider, reference current
        sources, CMFB path, and compensation capacitor — was implemented in Virtuoso.
  \item Temperature sweep simulation from $-20\,^\circ$C to $+100\,^\circ$C was performed
        and the $g_m$ characteristics were extracted and analysed.
\end{itemize}

\subsection*{What Was Not Achieved and Why}
\begin{itemize}[leftmargin=1.5em]
  \item The $< 1\%$ $g_m$ variation specification was \textbf{not met}. The simulated
        device-level $g_m$ varies by $\approx 19\%$ across the temperature range.
  \item The root causes identified are: (a) a mismatch between $I_{\text{ref}}$ and
        $2\Delta V$ that sets the wrong locking point; (b) the servo loop may not be
        actively regulating the bias — evidenced by the near-constant $V_{cm\_fb}$
        across temperature; (c) the probed quantity ($g_{m,\text{NM0}}$) is the
        device-level parameter, not the closed-loop $G_m$ regulated by the servo.
  \item These are \textbf{correctable issues} and do not indicate a fundamental flaw
        in the circuit topology, which is consistent with the reference design.
\end{itemize}

\subsection*{Learning Outcomes}
The project provided hands-on experience in:
hierarchical schematic design in Cadence Virtuoso,
transistor-level operating point analysis using Spectre,
understanding PVT sensitivity of $g_m$ and the role of negative feedback in suppressing it,
and systematic debugging of analog feedback loops — including identifying
loop connectivity issues, gain-bandwidth tradeoffs, and measurement methodology errors.

The fundamental circuit topology is sound and consistent with the reference design.
With further tuning of the resistive divider, a bandgap-referenced $I_{\text{ref}}$,
and loop stability verification via AC analysis, the circuit is expected to achieve
the $< 1\%$ regulation target.

% ════════════════════════════════════════════════════════════════════════════
\section*{References}
\addcontentsline{toc}{section}{References}

\begin{enumerate}[leftmargin=2em]
  \item B. Razavi, \textit{Design of Analog CMOS Integrated Circuits}, 2nd ed.,
        McGraw-Hill, 2016.
  \item F. Silveira, D. Flandre, and P. G. A. Jespers, ``A $g_m/I_D$ based methodology
        for the design of CMOS analog circuits and its application to the synthesis of a
        silicon-on-insulator micropower OTA,'' \textit{IEEE J. Solid-State Circuits},
        vol.~31, no.~9, pp.~1314--1319, Sep.~1996.
  \item P. R. Kinget, ``Device mismatch and tradeoffs in the design of analog circuits,''
        \textit{IEEE J. Solid-State Circuits}, vol.~40, no.~6, pp.~1212--1224, Jun.~2005.
  \item Reference design: \textit{Low-Transistor-Count Constant-$g_m$ Circuits}
        (as provided in project brief).
\end{enumerate}

\end{document}
